\documentclass[10pt,a4paper]{article}
\usepackage[margin=15mm]{geometry}
\usepackage{fontspec}
\setmainfont{DejaVuSerif.ttf}[Path=/usr/share/fonts/truetype/dejavu/,BoldFont=DejaVuSerif-Bold.ttf,ItalicFont=DejaVuSerif-Italic.ttf]
\setsansfont{DejaVuSans.ttf}[Path=/usr/share/fonts/truetype/dejavu/,BoldFont=DejaVuSans-Bold.ttf]
\setmonofont{DejaVuSansMono.ttf}[Path=/usr/share/fonts/truetype/dejavu/]
\usepackage{polyglossia}
\setmainlanguage{russian}
\usepackage{booktabs,array,tabularx}
\usepackage[unicode,hidelinks,pdfauthor={},pdftitle={ЛР4. Один агент и две MAS}]{hyperref}
\setlength{\parindent}{0pt}
\setlength{\parskip}{3pt}
\setlength{\tabcolsep}{5pt}
\pagestyle{plain}
\newcommand{\sectionline}[1]{\par\vspace{4pt}{\sffamily\bfseries #1}\par\vspace{2pt}}
\newcommand{\id}[1]{{\fontsize{8}{10}\selectfont\texttt{\detokenize{#1}}}}
\begin{document}
{\sffamily\bfseries\Large ЛР4. Один агент и две MAS на исправлении кода}

\sectionline{Гипотеза, реализация и условия}
Гипотеза: ревью и координация могут помогать исправлению программ, но при общем
лимите расход на сообщения и повторную историю способен перекрыть пользу.
Один и тот же экземпляр Qwen3.5-4B обслуживает все роли последовательно:

\begin{center}\small
\begin{tabular}{rl}
\textbf{single:}& одна роль $\leftrightarrow$ инструменты $\to$ независимая проверка;\\
\textbf{fixed:}& programmer $\to$ reviewer $\to$ programmer $\to$ проверка;\\
\textbf{coordinator:}& coordinator $\leftrightarrow$ \{developer, tester\} $\to$ проверка.
\\\textbf{Общий цикл:}& LLM $\to$ JSON $\to$ Engine $\to$ инструмент $\to$ наблюдение $\to$ LLM.
\end{tabular}
\end{center}
Роли имеют отдельные истории и общий артефакт. Все обращения к LLM и инструментам
идут через Engine; успех требует \texttt{engine.finish()}. Fixed всегда выполняет
три этапа (лимиты 10/6/12 действий), явно передаёт ревью. Координатор валидирует
точную JSON-схему и роль; до 10 решений и 3 действий исполнителя на делегирование.
Общие пределы: 48 вызовов модели, 80 инструментальных действий.

Сервер: vLLM на RTX 5070, BF16, окно 65\,536, thinking выключен,
temperature=0, seed=41729, send\_seed=false. Точный API ID:\par
\id{/home/maria/Projects/IAD_MAG_neuroslop/lab02_students/Qwen3.5-4B}.\par
Endpoint: \id{http://192.168.3.42:8000/v1}; клиент: Python 3.11.5, Linux/WSL2.
Выход/время эпизода: T --- 1800 токенов/300 с; C --- 4000/600 с;
M --- 6000/900 с. Цены не заданы, денежная стоимость неизвестна.

\sectionline{Полный итоговый эксперимент: 696 эпизодов}
T: $50\times3\times2$ бюджета $\times2$ повтора = 600;
C: $12\times3\times2=72$; M: $4\times3\times2=24$.
Бюджеты T выбраны на dev; C/M оставлены 60\,000 после технических пилотов.
Перед test созданы три freeze. Все 696 результатов проверены по manifest,
хешам артефактов, суммам usage и расходу ролей; пропусков нет.

\begin{center}\fontsize{8.6}{10.5}\selectfont
\begin{tabular}{llrrrrrr}
\toprule
Блок/бюджет & Схема & Успех/N & Токены, тыс. & Время, с & Вызовы & Бюджет исч. & Usage ?\\
\midrule
T/24k & single & 43/100 & 24.35 & 9.53 & 6 & 57 & 0\\
 & fixed & 14/100 & 25.94 & 16.63 & 7 & 86 & 0\\
 & coordinator & 26/100 & 25.41 & 17.31 & 7 & 73 & 1\\
T/48k & single & 69/100 & 25.39 & 12.55 & 6 & 29 & 2\\
 & fixed & 61/100 & 32.90 & 17.90 & 8 & 16 & 0\\
 & coordinator & 68/100 & 32.42 & 19.54 & 9 & 32 & 0\\
\midrule
C/60k & single & 2/24 & 65.55 & 83.40 & 7 & 18 & 4\\
 & fixed & 3/24 & 63.92 & 85.24 & 7 & 15 & 3\\
 & coordinator & 4/24 & 64.95 & 96.51 & 7 & 12 & 8\\
\midrule
M01--03/60k & single & 0/6 & 65.16 & 56.14 & 4 & 6 & 0\\
 & fixed & 0/6 & 72.64 & 266.43 & 5 & 6 & 0\\
 & coordinator & 0/6 & 69.05 & 99.06 & 5 & 6 & 0\\
\bottomrule
\end{tabular}
\end{center}
{\footnotesize Токены, время и вызовы --- медианы. Для токенов учтены $N-$Usage?
эпизодов, для времени --- все, кроме одного прерванного T/coordinator/24k
($n=99$); для вызовов там также $n=99$. Остальные знаменатели полные.
«Успех» означает завершённый процесс и прохождение открытых тестов;
закрытая правильность C/M не установлена, у M \texttt{success=null}.
Полные статусы и доли исчерпания --- \id{analysis/final/summary.csv}.}

\textbf{M04 отдельно (по 2 эпизода на схему):} у всех формат 0/2, открытый пример
не проверен (программа не сохранена), закрытый результат неизвестен.
Медианы токенов/секунд/вызовов: single 105\,077/23.35/4;
fixed 105\,100/21.61/4; coordinator 68\,530/30.99/4.
Все 6 эпизодов исчерпали бюджет; готовых ответов для сравнения нет.

\newpage
{\sffamily\bfseries\large Парное сравнение, трассы и ограничения}

\sectionline{Сравнение на одинаковых task/budget/repeat}
В строке по 100 пар. «Только MAS» --- победа команды; ошибки API, прерывание и
исчерпание бюджета остаются неуспехами. Оба повтора одной задачи зависимы.
\begin{center}\small
\begin{tabular}{llrrrr}
\toprule
Бюджет & MAS & Оба успешны & Только MAS & Только single & Оба неуспешны\\
\midrule
24k & fixed & 14 & 0 & 29 & 57\\
24k & coordinator & 20 & 6 & 23 & 51\\
48k & fixed & 54 & 7 & 15 & 24\\
48k & coordinator & 56 & 12 & 13 & 19\\
\bottomrule
\end{tabular}
\end{center}
При 48k разница с single: fixed $-8$ п.п., coordinator $-1$ п.п.; близкие общие
доли скрывают разные решённые задачи. В C у fixed одна победа и ноль поражений,
у coordinator --- две и ноль (по 24 пары; таймауты сохранены).
В теме «Ресурсы и ограничения» на T/24k coordinator успешен 10/20 против 6/20
у single; это локальное наблюдение, не общее преимущество MAS.

Совпадение фактических результатов открытых тестов (не текста кода): fixed
62/100 и 72/98 доступных пар при 24k/48k, coordinator --- 66/97 и 73/98.
В C --- 19/24 у каждой команды. Остальные сравнимые пары различаются;
недоступные наблюдения не считаются совпадением.
Медиана парной разности токенов MAS$-$single на 48k: fixed $+7456$,
coordinator $+3559$ (по 98 пар с полным usage); времени $+5.51$ и $+4.47$ с
(по 100 пар). Это оценки доступных пар, не доказательство разницы стоимости
всей выборки. Среди двух успешных процессов разности токенов $+9057$
и $+4391$ соответственно (54 и 56 пар); это отобранное подмножество.

\sectionline{Три трассы: слова роли, инструмент, расход}
\textbf{1. Полезное ревью: \id{T43_fixed_b48000_r2}.}
Первый программист заявил об исправлении, но не применил правку: тесты дали
7/10. Reviewer указал замену суммы времён предшественников на максимум;
его попытка редактирования была запрещена. Замечание получил заключительный
программист, сохранил правку, получил 10/10 и независимое подтверждение.
Успех за 37\,696 токенов; reviewer --- 10\,932 (29\%). Single в той же паре
тоже дошёл до 10/10, но его finish пришёл уже после лимита: 55\,568 токенов,
неуспех. Польза здесь в последовательности работы и своевременном завершении,
а не в уникальной способности получить правильный код.

\textbf{2. Межмодульная ошибка: \id{M01_fixed_b60000_r1}.}
Первые ответы программиста были невалидным JSON либо неверной схемой
replace\_files. Только четвёртый вызов сохранил десять файлов, нарушив интерфейс:
\id{cancellation.py} импортирует \id{release}, отсутствующую в
\id{inventory.py}. Итоговая независимая проверка --- ImportError и 0/12.
Все 61\,217 токенов ушли первому программисту; reviewer не был вызван до
остановки. Это ошибка модельной правки, а не прежний отказ чтения stdlib.

\textbf{3. Корпус и история: \id{M04_single_b60000_r1}.}
Модель прочитала первые 100 строк shard\_01, затем запросила 500 (инструмент
разрешает 1--100), после отказа повторно прочитала те же 100 строк.
Вход вырос с 7327 до 45\,077 токенов на запрос. Всего учтено 104\,811 входных
и 266 выходных токенов; последний ответ отклонён за бюджетом. Расчёт программы
и ответ не сохранены, формат невалиден, открытый пример не проверен.
Большое окно само по себе не обеспечивает эффективного использования корпуса.

\sectionline{Изменения протокола, ограничения и вывод}
До test раскрываемо изменены: HTTP-доступ к LAN, общее уточнение finish/yield,
доступ sandbox к stdlib Anaconda и проверка этого импорта в doctor. Это отступления
от запретов исходного комплекта, принятые без отдельного согласования преподавателя;
код baseline и данные не менялись. Все пилоты до/после приложены; 63 offline-теста
прошли. Настройки после freeze не менялись.
В T сохранены 2 ошибки транспорта и 1 пользовательское прерывание, в C ---
15 таймаутов API. Полный usage неизвестен в 18 эпизодах, не заменён нулём.
Расход сверх лимита учитывается целиком, действие ответа не применяется.
Вся выборка из 696 эпизодов сохранена, выборочных повторов нет.

На этой модели и наборе single лучше при тесном бюджете; при 48k coordinator
почти равен ему по успеху, но имеет дополнительные расходы. M не завершён при
60k: нельзя обобщать это на все модели или утверждать закрытую неправильность.
Два повтора, одна модель, сетевые сбои и один интеграционный реестр ограничивают
выводы; программная обработка M04 не является чистым тестом окна контекста.
Полные таблицы, события трёх трасс и воспроизведение ---
\id{analysis/}, \id{runs/}, \id{SUBMISSION_README.md}.
\end{document}
